\begin{table*}[t]
\centering
\caption{Comparison with four traditional planning baselines on the held-out TC segments from the same 500-TC train/test split. Surface values are means over the 4/8/12/16-SUV fleets, both speed limits, and all five deployment radii. Subsurface values are means over all seven depths and both propulsion limits. Smaller errors and paths are better; higher current and coverage are better.}
\label{tab:classical_planners}
\small
\resizebox{\textwidth}{!}{%
\begin{tabular}{@{}lrrrr@{}}
\toprule
\multicolumn{5}{l}{\textbf{Panel A: surface tracking}}\\
Method & Final error (km) & Time-averaged distance (km) & \multicolumn{2}{c}{Path per SUV (km)}\\
\midrule
Parameter-sharing IPPO & \textbf{356.2} & \textbf{356.6} & \multicolumn{2}{c}{\textbf{462.3}} \\
Potential field & 413.4 & 404.2 & \multicolumn{2}{c}{728.8} \\
Receding-horizon A* & 457.2 & 429.4 & \multicolumn{2}{c}{823.9} \\
Lawnmower coverage & 554.7 & 478.4 & \multicolumn{2}{c}{509.1} \\
Information greedy & 466.2 & 429.4 & \multicolumn{2}{c}{528.7} \\
\midrule
\multicolumn{5}{l}{\textbf{Panel B: subsurface observation}}\\
Method & Mean current (m/s) & Nominal coverage (km$^2$) & Path (km) & Current-weighted coverage ($10^3$ m\,s$^{-1}$\,km$^2$)\\
\midrule
Parameter-sharing IPPO & 0.408 & \textbf{3506} & 123.0 & \textbf{1.599} \\
Potential field & \textbf{0.452} & 2048 & 107.3 & 1.050 \\
Receding-horizon A* & 0.427 & 2310 & \textbf{95.8} & 1.097 \\
Lawnmower coverage & 0.408 & 2914 & 110.8 & 1.286 \\
Information greedy & 0.429 & 1549 & 98.8 & 0.739 \\
\bottomrule
\end{tabular}
}
\end{table*}
